\SourcePage{245}
\SGAsection{5}{诺特的 $J$ 进与 AR-$J$ 进投射系}
本节始终假定给定一个交换\Term{noetherian}{诺特}环 $A$、$A$ 的一个\Term{ideal}{理想} $J$，以及一个\Term{abelian-a-category}{阿贝尔 $A$ 范畴} $C$。

\textbf{5.1. 一般性质。}
\begin{SGAresult}{定义}{5.1.1}
若\Term{category}{范畴} $\operatorname{Hom}(\mathbb N^\circ,C)$ 的\Term{object}{对象} $X$ 是\Term{j-adic-system}{$J$ 进}的，且其各\Term{set-component}{分量}都是\Term{category}{范畴} $C$ 的\Term{noetherian}{诺特}\Term{object}{对象}（相应地，\Term{artinian}{阿廷}\Term{object}{对象}），则称 $X$ 为\Term{noetherian}{诺特}（相应地，\Term{artinian}{阿廷}）的\Term{j-adic-system}{$J$ 进投射系}。
\end{SGAresult}

以 $J\text{-}\mathrm{adn}(C)$（相应地，$J\text{-}\mathrm{adf}(C)$）表示 $P=\operatorname{Hom}(\mathbb N^\circ,C)$ 中由\Term{noetherian}{诺特}（相应地，\Term{artinian}{阿廷}）的\Term{j-adic-system}{$J$ 进投射系}构成的\Term{full-subcategory}{满子范畴}。显然，由 3.1.3，这类\Term{projective-system}{投射系}的\Term{extension}{扩张}仍属于同一类型。

\begin{SGAresult}{命题}{5.1.2}
设 $X$ 是 $P$ 的一个\Term{j-adic-system}{$J$ 进}\Term{object}{对象}。$X$ 为\Term{noetherian}{诺特}（相应地，\Term{artinian}{阿廷}）的\Term{j-adic-system}{$J$ 进投射系}，当且仅当它是\Term{j-adic-system}{$J$ 进}的，且 $X_0$ 是\Term{noetherian}{诺特}（相应地，\Term{artinian}{阿廷}）\Term{object}{对象}。
\end{SGAresult}
\begin{proof}
设 $X$ 是一个\Term{j-adic-system}{$J$ 进投射系}，而 $X_0$ 是\Term{noetherian}{诺特}（相应地，\Term{artinian}{阿廷}）的。考虑\Term{canonical}{典范}\Term{epimorphism}{满态射}（4.2.4）
\[
\operatorname{gr}_J(A)\otimes_A X_0\longrightarrow\operatorname{grs}(X).
\]
由于 $X_0$ 是\Term{noetherian}{诺特}（相应地，\Term{artinian}{阿廷}）的，且对每个 $n$，$J^n$ 都是\Term{finite-type}{有限型}的，这个\Term{epimorphism}{满态射}的存在表明，对每个整数 $n$，\Term{canonical-morphism}{典范态射}
\[
X_n\longrightarrow X_{n-1}
\]
的核都是\Term{noetherian}{诺特}（相应地，\Term{artinian}{阿廷}）的。由此对 $n$ 归纳，立即得到各 $X_n$ 都是\Term{noetherian}{诺特}（相应地，\Term{artinian}{阿廷}）的。
\end{proof}

\SourcePage{246}
\begin{SGAresult}{定义}{5.1.3}
设 $X$ 是一个以 $\mathbb N$ 为指标、取值于 $C$ 的\Term{projective-system}{投射系}。若它是\Term{ar-j-adic-system}{AR-$J$ 进}的，且各\Term{set-component}{分量}都是\Term{noetherian}{诺特}（相应地，\Term{artinian}{阿廷}）\Term{object}{对象}，则称它为\Term{noetherian}{诺特}（相应地，\Term{artinian}{阿廷}）的\Term{ar-j-adic-system}{AR-$J$ 进投射系}。
\end{SGAresult}
以 $\mathrm{AR}\text{-}J\text{-}\mathrm{adn}(C)$（相应地，$\mathrm{AR}\text{-}J\text{-}\mathrm{adf}(C)$）表示 $\operatorname{Hom}_{\mathrm{AR}}(\mathbb N^\circ,C)$（2.4）中，由\Term{noetherian}{诺特}（相应地，\Term{artinian}{阿廷}）的\Term{ar-j-adic-system}{AR-$J$ 进投射系}的像生成的\Term{full-subcategory}{满子范畴}。
于是显然，\Term{restriction}{限制} $P_{\mathrm{AR}}$（2.4）所得的\Term{canonical-functor}{典范函子}
\[
P_{\mathrm{AR}}^n:J\text{-}\mathrm{adn}(C)\longrightarrow\mathrm{AR}\text{-}J\text{-}\mathrm{adn}(C),
\]
\[
P_{\mathrm{AR}}^f:J\text{-}\mathrm{adf}(C)\longrightarrow\mathrm{AR}\text{-}J\text{-}\mathrm{adf}(C)
\]
都是\Term{category-equivalence}{范畴等价}。

简要回顾 Hilbert \Term{theorem}{定理}及其证明。
\begin{SGAresult}{\Term{theorem}{定理}}{5.1.4}
\textbf{（Hilbert \Term{theorem}{定理}。）} 设 $A$ 是一个\Term{commutative-ring}{交换环}，$C$ 是一个\Term{abelian-a-category}{阿贝尔 $A$ 范畴}，$M$ 是 $C$ 的一个\Term{finite-type}{有限型}（相应地，\Term{noetherian}{诺特}）\Term{object}{对象}。

对于每个以 $\mathbb N$ \Term{graded}{分次}的\Term{finite-type}{有限型} $A$ \Term{algebra}{代数} $B$，\Term{graded-module}{分次 $(B,C)$ 模} $B\otimes_A M$ 都是\Term{finite-type}{有限型}（相应地，\Term{noetherian}{诺特}）的。

（注意：若 $B=\bigoplus_{n\geq0}B_n$，则 $B\otimes_A M$ 表示\Term{graded-object}{分次对象} $(B_n\otimes_A M)_{n\in\mathbb N}$。）
\end{SGAresult}
\begin{proof}
立即可归结为 $B$ 是带有通常\Term{graded}{分次}的\Term{polynomial-ring}{多项式环} $A[T]$ 的情形。此时\Term{graded-module}{分次 $(B,C)$ 模}
\[
A[T]\otimes_A M=(M_n)_{n\in\mathbb N}
\]
由如下方式得到：对每个 $n$ 置 $M_n=M$，而 $T$ 对每个整数 $n$ 都以 $M$ 的\Term{identity-morphism}{恒等态射}把 $M_n$ 映到 $M_{n+1}$。

设 $N$ 是 $A[T]\otimes_A M$ 的一个\Term{graded-module}{分次 $(A[T],C)$ 子模}。$N$ 的各\Term{set-component}{分量} $(N_n)_{n\in\mathbb N}$ 给出 $M$ 的\Term{subobject}{子对象}的递增列
\[
N_0\subset N_1\subset\cdots\subset N_n\subset\cdots.
\]
\SourcePage{247}
现在假设 $M$ 是\Term{noetherian}{诺特}（相应地，\Term{finite-type}{有限型}）的，并设
\[
N^1\subset N^2\subset\cdots\subset N^p\subset\cdots
\]
是 $A[T]\otimes_A M$ 的\Term{graded-module}{分次 $(A[T],C)$ 子模}的一个递增列；在第一种假设下（即 $M$ 为\Term{finite-type}{有限型}时），还假定其并为 $A[T]\otimes_A M$。
考虑如下双指标表：
\[
\begin{tikzcd}[row sep=1.6em,column sep=2em]
N_0^1\arrow[r,phantom,"\subset"]\arrow[d,phantom,"\cap" description]&N_0^2\arrow[r,phantom,"\subset\ \cdots\ \subset"]\arrow[d,phantom,"\cap" description]&N_0^p\arrow[r,phantom,"\subset\ \cdots"]\arrow[d,phantom,"\cap" description]&{}\\
N_1^1\arrow[d,phantom,"\vdots" description]&N_1^2\arrow[d,phantom,"\vdots" description]&N_1^p\arrow[d,phantom,"\vdots" description]&{}\\
N_n^1\arrow[r,phantom,"\subset"]\arrow[d,phantom,"\cap" description]&N_n^2\arrow[r,phantom,"\subset\ \cdots\ \subset"]\arrow[d,phantom,"\cap" description]&N_n^p\arrow[r,phantom,"\subset\ \cdots"]\arrow[d,phantom,"\cap" description]&{}\\
\vdots&\vdots&\vdots&{}
\end{tikzcd}
\]
由关于 $M$ 的假设，显然存在一对 $(p_0,n_0)$，使得 $N_{n_0}^{p_0}$ 极大（相应地，等于 $M$）。取整数 $p_1\geq p_0$，使 $N_n^{p_1}$（$n\leq n_0$）是各列 $(N_n^p)_{p\in\mathbb N}$ 的最大元素，则显然，只要 $p\geq p_1$ 就有 $N^p=N^{p_1}$。这便在两种情形下都给出了结论。
\end{proof}

\begin{SGAresult}{推论}{5.1.5}
再假定\Term{category}{范畴} $C$ 具有可数\Term{direct-sum}{直和}。则对 $C$ 的每个\Term{finite-type}{有限型}（相应地，\Term{noetherian}{诺特}）\Term{object}{对象} $M$ 及每个\Term{finite-type}{有限型} $A$ \Term{algebra}{代数} $B$，$(B,C)$ 模 $B\otimes_A M$ 都是\Term{finite-type}{有限型}（相应地，\Term{noetherian}{诺特}）的。
\end{SGAresult}
\begin{proof}
归结为 $B=A[T]$ 的情形。设 $N$ 是 $A[T]\otimes_A M$ 的一个 $(A[T],C)$ \Term{submodule}{子模}，此处将后者看作 $C$ 的\Term{object}{对象}。以 $\operatorname{gr}(N)$ 表示 $N$ 相对于由 $A[T]\otimes_A M$ 的自然\Term{graded-object}{分次对象}\Term{filtration}{滤过}诱导的\Term{filtration}{滤过}所得的\Term{associated-graded}{伴随分次对象}。$\operatorname{gr}(N)$ 是\Term{graded-module}{分次 $(A[T],C)$ 模} $A[T]\otimes_A M$ 的一个\Term{graded-module}{分次 $(A[T],C)$ 子模}。
\SourcePage{248}
现在设
\[
N^0\subset N^1\subset\cdots\subset N^p\subset\cdots
\]
是 $A[T]\otimes_A M$ 的 $(A[T],C)$ \Term{submodule}{子模}的一个递增列，并在第一种假设下假定其并为 $A[T]\otimes_A M$。在第一种假设下（$M$ 为\Term{finite-type}{有限型}），存在整数 $i$，使 $\operatorname{gr}(N^i)=A[T]\otimes_A M$，从而 $N^i=A[T]\otimes_A M$。在第二种假设下（$M$ 为\Term{noetherian}{诺特}），$\operatorname{gr}(N^i)$ 构成的列最终\Term{invariant}{不变}，故 $N^i$ 构成的列也最终\Term{invariant}{不变}。
\end{proof}
以下主要通过下述命题使用 Hilbert \Term{theorem}{定理}。

\begin{SGAresult}{命题}{5.1.6}
设 $A$ 是一个交换\Term{noetherian}{诺特}环，$J$ 是 $A$ 的一个\Term{ideal}{理想}，$C$ 是一个\Term{abelian-a-category}{阿贝尔 $A$ 范畴}。对每个\Term{noetherian}{诺特}的\Term{j-adic-system}{$J$ 进投射系} $X$（5.1.1），其\Term{associated-graded}{分次对象} $\operatorname{grs}(X)$ 都是一个\Term{noetherian}{诺特}的 $(\operatorname{gr}_J(A),C)$ 模。
\end{SGAresult}
\begin{proof}
由于 $\operatorname{gr}_J(A)$ 是\Term{finite-type}{有限型} $A$ \Term{algebra}{代数}，Hilbert \Term{theorem}{定理}（5.1.4）表明 $\operatorname{gr}_J(A)\otimes_A X_0$ 是一个\Term{noetherian}{诺特}的\Term{graded-module}{分次 $(\operatorname{gr}_J(A),C)$ 模}。由\Term{epimorphism}{满态射}（4.2.4）
\[
\operatorname{gr}_J(A)\otimes_A X_0\longrightarrow\operatorname{grs}(X)
\]
的存在，即得所述断言。
\end{proof}

\textbf{5.2. $J\text{-}\mathrm{adn}(C)$ 与 $\mathrm{AR}\text{-}J\text{-}\mathrm{adn}(C)$ 的结构。}
本小节以 $\mathscr E_C$，或在不致混淆时以 $\mathscr E$，表示 $\operatorname{Hom}(\mathbb N^\circ,C)$ 中由\Term{noetherian}{诺特}的\Term{ar-j-adic-system}{AR-$J$ 进投射系}构成的\Term{full-subcategory}{满子范畴}。

\begin{SGAresult}{命题}{5.2.1}
\Term{category}{范畴} $\mathscr E_C$ 在 $\operatorname{Hom}(\mathbb N^\circ,C)$ 中对取核与\Term{cokernel}{余核}\Term{stable}{稳定}。换言之，$\mathscr E_C$ 是一个\Term{abelian-category}{阿贝尔范畴}，且\Term{inclusion-functor}{包含函子}
\SourcePage{249}
\[
\mathscr E_C\longrightarrow\operatorname{Hom}(\mathbb N^\circ,C)
\]
是\Term{exact}{正合}的。
\end{SGAresult}
\begin{proof}
设 $u:X\to Y$ 是 $\mathscr E_C$ 的一个\Term{morphism}{态射}。我们依次证明 $u$ 的\Term{cokernel}{余核}、像与核都是\Term{noetherian}{诺特}的\Term{ar-j-adic-system}{AR-$J$ 进投射系}。（当然，关于 $u$ 的像的这一断言，在证明中仅起技术性中间步骤的\Term{action}{作用}。）

(a) $\operatorname{Coker}(u)$ 是\Term{noetherian}{诺特}的\Term{ar-j-adic-system}{AR-$J$ 进投射系}。由于 $X$ 满足\Term{mlar-condition}{$(\mathrm{MLAR})$}，而 $Y$ 是\Term{ar-j-adic-system}{AR-$J$ 进}的，由 3.2.4 (i)，$\operatorname{Coker}(u)$ 是\Term{ar-j-adic-system}{AR-$J$ 进}的。

(b) $\operatorname{Im}(u)$ 是\Term{noetherian}{诺特}的\Term{ar-j-adic-system}{AR-$J$ 进投射系}。利用\Term{exact-sequence}{正合列}
\[
0\longrightarrow\operatorname{Im}(u)\longrightarrow Y\longrightarrow\operatorname{Coker}(u)\longrightarrow0,
\]
这一断言将由下面的\Term{lemma}{引理} 5.2.1.1 得出。
\end{proof}

\begin{SGAresult}{\Term{lemma}{引理}}{5.2.1.1}
设
\[
0\longrightarrow L\longrightarrow M\longrightarrow N\longrightarrow0
\]
是\Term{category}{范畴} $\operatorname{Hom}(\mathbb N^\circ,C)$ 中的一个\Term{exact-sequence}{正合列}。若 $M$ 与 $N$ 都是\Term{noetherian}{诺特}的\Term{ar-j-adic-system}{AR-$J$ 进投射系}，则 $L$ 也是如此。
\end{SGAresult}
\begin{proof}
显然，只须证明它是\Term{ar-j-adic-system}{AR-$J$ 进}的。与 3.2.4 的证明一样，先可归结到 $M$ 与 $N$ 都是\Term{strict-projective-system}{严格}的情形，再进一步归结到它们都是\Term{noetherian}{诺特}的\Term{j-adic-system}{$J$ 进投射系}的情形。这时由 3.1.3 (ii)，$L$ 是\Term{strict-projective-system}{严格}的。\Term{projective-system}{投射系} $L$ 是\Term{strict-projective-system}{严格}的、\Term{j-adapted}{适合于 $J$}的，而且它的\Term{associated-graded}{分次对象}作为 $\operatorname{grs}(M)$ 的\Term{graded-object}{分次}\Term{subobject}{子对象}，是\Term{noetherian}{诺特}的：$\operatorname{grs}(M)$ 是 $\operatorname{gr}_J(A)\otimes_A M_0$ 的商，因而是\Term{noetherian}{诺特}的（5.1.6）。由 Artin--Rees \Term{lemma}{引理}（4.2.6）即得断言。
\end{proof}
\begin{proof}[命题 5.2.1 的证明（续）]
(c) $\operatorname{Ker}(u)$ 是\Term{noetherian}{诺特}的\Term{ar-j-adic-system}{AR-$J$ 进投射系}。只须将 5.2.1.1 应用于下面显然的\Term{exact-sequence}{正合列}：
\[
0\longrightarrow\operatorname{Ker}(u)\longrightarrow X\longrightarrow\operatorname{Im}(u)\longrightarrow0.
\]
\end{proof}

\SourcePage{250}
\begin{SGAresult}{命题}{5.2.2}
\Term{abelian-category}{阿贝尔范畴} $\mathrm{AR}\text{-}J\text{-}\mathrm{adn}(C)$ 的各\Term{object}{对象}都是\Term{noetherian}{诺特}的。
\end{SGAresult}
\begin{proof}
立即可归结到证明以下事实。设给定一个\Term{noetherian}{诺特}的\Term{j-adic-system}{$J$ 进}\Term{object}{对象} $X$、一列\Term{noetherian}{诺特}的\Term{j-adic-system}{$J$ 进}\Term{object}{对象} $(Y^n)_{n\in\mathbb N}$，以及 $\operatorname{Hom}(\mathbb N^\circ,C)$ 中的\Term{morphism}{态射}
\[
v_n:Y^n\longrightarrow X,
\]
满足：

(i) $v_n$ 在 $\mathrm{AR}\text{-}J\text{-}\mathrm{adn}(C)$ 中变成\Term{monomorphism}{单态射}；

(ii) 对每个整数 $n$，$v_n$ 在 $\mathrm{AR}\text{-}J\text{-}\mathrm{adn}(C)$ 中都通过 $v_{n+1}$ \Term{resolution}{分解}。

则当 $n$ 充分大时，$v_n$ 在 $\mathrm{AR}\text{-}J\text{-}\mathrm{adn}(C)$ 中在\Term{isomorphism}{同构}意义下最终\Term{invariant}{不变}。

而 $v_n$ 在 $\mathrm{AR}\text{-}J\text{-}\mathrm{adn}(C)$ 中通过 $v_{n+1}$ \Term{resolution}{分解}，意味着存在一个整数 $r$ 和一个\Term{morphism}{态射}
\[
f:Y^n_{[r]}\longrightarrow Y^{n+1},
\]
使得若以 $a$ 表示\Term{canonical-morphism}{典范态射}
\[
Y^n_{[r]}\longrightarrow Y^n,
\]
则有
\[
v_n\circ a=v_{n+1}\circ f.
\]
由于\Term{morphism}{态射} $a$ 是\Term{strict-morphism}{严格}的，特别可得 $v_n$ 的像包含于 $v_{n+1}$ 的像。以 $\operatorname{Im}(v_n)$ 代替各 $Y^n$，便归结为证明：$X$ 的任一\Term{strict-projective-system}{严格}子\Term{projective-system}{投射系}的递增列都最终\Term{invariant}{不变}。设 $(Z^n)_{n\in\mathbb N}$ 是这样一列；由于 $\operatorname{grs}(X)$ 是\Term{noetherian}{诺特}的，$\operatorname{grs}(Z^n)$ 构成的递增列最终\Term{invariant}{不变}，故 $Z^n$ 构成的列也最终\Term{invariant}{不变}。
\end{proof}
综上得到：

\SourcePage{251}
\begin{SGAresult}{\Term{theorem}{定理}}{5.2.3}
\Term{category}{范畴} $J\text{-}\mathrm{adn}(C)$ 与 $\mathrm{AR}\text{-}J\text{-}\mathrm{adn}(C)$ 都是\Term{abelian-category}{阿贝尔}且\Term{noetherian}{诺特}的。
\end{SGAresult}
\begin{proof}
对 $\mathrm{AR}\text{-}J\text{-}\mathrm{adn}(C)$，这正是刚才所得的结论；而 $J\text{-}\mathrm{adn}(C)$ 与它\Term{category-equivalence}{等价}。
\end{proof}

\begin{SGAresult}{命题}{5.2.4}
\textbf{（对\Term{group-extension}{扩张}的\Term{stable}{稳定}性。）} 设
\[
0\longrightarrow X\longrightarrow Y\longrightarrow Z\longrightarrow0
\]
是 $\operatorname{Hom}(\mathbb N^\circ,C)$ 中的一个\Term{exact-sequence}{正合列}，且对每个 $n$ 都有 $J^{n+1}Y_n=0$。假定：

(i) $Z$ 是\Term{ar-j-adic-system}{AR-$J$ 进}的；

(ii) $X$ 是\Term{artinian}{阿廷}的\Term{ar-j-adic-system}{AR-$J$ 进投射系}。

则 $Y$ 是\Term{ar-j-adic-system}{AR-$J$ 进}的。
\end{SGAresult}
\begin{proof}
与 3.2.4 的证明一样，归结为 $Y$ 与 $Z$ 都是\Term{strict-projective-system}{严格}的、且 $Z$ 是\Term{j-adic-system}{$J$ 进}的情形。这时由 3.1.3 (ii)，$X$ 是\Term{strict-projective-system}{严格}的。于是存在整数 $r$，使\Term{projective-system}{投射系}
\[
\bigl(X_{n+r}/J^{n+1}X_{n+r}\bigr)_{n\in\mathbb N}
\]
是\Term{artinian}{阿廷}的\Term{j-adic-system}{$J$ 进投射系}。由于 $X$ 是它的一个商，且与它\Term{ar-isomorphism}{AR 同构}，只须证明下述\Term{lemma}{引理}。
\end{proof}

\begin{SGAresult}{\Term{lemma}{引理}}{5.2.4.1}
设
\[
0\longrightarrow T\longrightarrow X\longrightarrow Y\longrightarrow Z\longrightarrow0
\]
是 $\operatorname{Hom}(\mathbb N^\circ,C)$ 中的一个\Term{exact-sequence}{正合列}，且对每个 $n$ 都有 $J^{n+1}Y_n=0$。若 $X$ 是\Term{artinian}{阿廷}的\Term{j-adic-system}{$J$ 进投射系}，$Z$ 是\Term{j-adic-system}{$J$ 进}的，而 $T$ 是\Term{ar-null}{AR 零}的，则 $Y$ 是\Term{ar-j-adic-system}{AR-$J$ 进}的。
\end{SGAresult}
\begin{proof}
为证明此\Term{lemma}{引理}，先证明对每个整数 $n$，显然的\Term{projective-system}{投射系}
\[
\bigl(Y_{n+k}/J^{n+1}Y_{n+k}\bigr)_{k\in\mathbb N}
\]
都是\Term{essentially-constant}{本质常值}的。以 $S_n^k$ 表示显然\Term{morphism}{态射}
\[
X_{n+k}/J^{n+1}X_{n+k}\longrightarrow Y_{n+k}/J^{n+1}Y_{n+k}
\]
的核。
\SourcePage{252}
考虑下面显然的\Term{exact}{正合}\Term{commutative-diagram}{交换图}：
\[
\resizebox{\textwidth}{!}{\begin{tikzcd}[column sep=small,ampersand replacement=\&]
0\arrow[r]\&S_n^k\arrow[r]\arrow[d]\&X_{n+k}/J^{n+1}X_{n+k}\arrow[r]\arrow[d,"\sim"']\&Y_{n+k}/J^{n+1}Y_{n+k}\arrow[r]\arrow[d]\&Z_{n+k}/J^{n+1}Z_{n+k}\arrow[r]\arrow[d,"\sim"]\&0\\
0\arrow[r]\&T_n\arrow[r]\&X_n\arrow[r]\&Y_n\arrow[r]\&Z_n\arrow[r]\&0.
\end{tikzcd}}
\]
固定整数 $n$ 后，各\Term{object}{对象} $S_n^k$ 可视为 $T_n$ 的\Term{subobject}{子对象}，且构成一个递减列。由于 $T_n$ 是\Term{artinian}{阿廷}的，存在整数 $k_0$，使得对所有 $k\geq k_0$ 都有 $S_n^k=S_n^{k_0}$。由此立即得到，\Term{canonical-morphism}{典范态射}
\[
Y_{n+k}/J^{n+1}Y_{n+k}\longrightarrow Y_{n+k_0}/J^{n+1}Y_{n+k_0}
\]
都是\Term{isomorphism}{同构}。

现在以 $Y'$ 表示如下定义的\Term{projective-system}{投射系}：在\Term{object}{对象}上置
\[
Y'_n=\varprojlim_r Y_{n+r}/J^{n+1}Y_{n+r},
\]
而\Term{morphism}{态射}按显然的方式定义。\Term{projective-system}{投射系} $Y'$ 是\Term{j-adic-system}{$J$ 进}的。此外，“取\Term{limit}{极限}”得到如下\Term{form}{形式}的\Term{exact}{正合}\Term{commutative-diagram}{交换图}：
\[
\begin{tikzcd}[column sep=1.5em]
0\arrow[r]&S\arrow[r]\arrow[d]&X\arrow[r]\arrow[d,"\mathrm{id}"']&Y'\arrow[r]\arrow[d]&Z\arrow[r]\arrow[d,"\mathrm{id}"]&0\\
0\arrow[r]&T\arrow[r]&X\arrow[r]&Y\arrow[r]&Z\arrow[r]&0.
\end{tikzcd}
\]
由于 $T$ 是\Term{ar-null}{AR 零}的，对此图应用\Term{snake-lemma}{蛇引理}可知，显然\Term{morphism}{态射} $Y'\to Y$ 是一个\Term{ar-isomorphism}{AR 同构}，所以 $Y$ 是\Term{ar-j-adic-system}{AR-$J$ 进}的。
\end{proof}

\SourcePage{253}
\begin{SGAresult}{推论}{5.2.5}
设
\[
0\longrightarrow X\longrightarrow Y\longrightarrow Z\longrightarrow0
\]
是 $\operatorname{Hom}(\mathbb N^\circ,C)$ 中的一个\Term{exact-sequence}{正合列}，且对每个 $n$ 都有 $J^{n+1}Y_n=0$。假定：

(i) $X$ 是\Term{artinian}{阿廷}的\Term{ar-j-adic-system}{AR-$J$ 进投射系}（相应地，既\Term{artinian}{阿廷}又\Term{noetherian}{诺特}）；

(ii) $Z$ 是\Term{artinian}{阿廷}（相应地，\Term{noetherian}{诺特}）的\Term{ar-j-adic-system}{AR-$J$ 进投射系}。

则 $Y$ 是\Term{artinian}{阿廷}（相应地，\Term{noetherian}{诺特}）的\Term{ar-j-adic-system}{AR-$J$ 进投射系}。
\end{SGAresult}

\textbf{5.3. \Term{noetherian}{诺特}的 AR-$J$ 进\Term{inverse-system}{投射系}与 $\partial$ \Term{functor}{函子}。}
设 $C$ 与 $D$ 是两个\Term{abelian-a-category}{阿贝尔 $A$ 范畴}，$F:C\to D$ 是一个\Term{additive-functor}{加性函子}。\Term{functor}{函子} $F$ 显然延拓为一个仍记作 $F$ 的\Term{functor}{函子}
\[
F:\operatorname{Hom}(\mathbb N^\circ,C)\longrightarrow\operatorname{Hom}(\mathbb N^\circ,D),
\]
它是\Term{additive}{加性}的，并且“与\Term{translation}{平移}交换”。

\begin{SGAresult}{命题}{5.3.1}
设 $C$ 与 $D$ 是两个\Term{abelian-a-category}{阿贝尔 $A$ 范畴}，而
\[
T=(T^i)_{i\in\mathbb N}
\]
是一个从 $C$ 到 $D$ 的 $A$ \Term{linear}{线性}\Term{exact}{正合}\Term{delta-functor}{$\partial$ 函子}。假定：

(i) 对 $C$ 的每个被 $J$ 的某个\Term{nilpotent}{幂零}化的\Term{noetherian}{诺特}\Term{object}{对象} $M$ 及每个整数 $n$，$T^n(M)$ 都是\Term{noetherian}{诺特}的。

(ii) 对每个被 $J$ \Term{annihilate}{零化}的\Term{finite-type}{有限型}\Term{graded-algebra}{分次 $A$ 代数} $B$，以及每个\Term{noetherian}{诺特}的\Term{graded-module}{分次 $(B,C)$ 模} $N$，各\Term{graded-module}{分次 $(B,D)$ 模} $T^n(N)$（$n\in\mathbb N$）都是\Term{noetherian}{诺特}的。

则对 $\operatorname{Hom}(\mathbb N^\circ,C)$ 的每个\Term{noetherian}{诺特}的\Term{ar-j-adic-system}{AR-$J$ 进}\Term{object}{对象} $X$，各 $T^n(X)$ 都是\Term{noetherian}{诺特}的\Term{ar-j-adic-system}{AR-$J$ 进投射系}。
\end{SGAresult}
\SourcePage{254}
\begin{proof}
由于\Term{functor}{函子} $T^n$ 与\Term{translation}{平移}交换，只须在 $X$ 为\Term{noetherian}{诺特}的\Term{j-adic-system}{$J$ 进投射系}时证明断言。这时 $\operatorname{grs}(X)$ 是一个\Term{noetherian}{诺特}的\Term{graded-module}{分次 $\operatorname{gr}_J(A)$ 模}（5.1.6），由假设 (ii) 及 Shih \Term{theorem}{定理}的一个推论的变式（参见附录 A.3），对每个整数 $n$ 都有：

(a) $T^n(X)$ 满足\Term{mlar-condition}{$(\mathrm{MLAR})$}；

(b) $\operatorname{grs}T^n(X)$ 是一个\Term{noetherian}{诺特}的\Term{graded-module}{分次 $(\operatorname{gr}_J(A),D)$ 模}。

由“Artin--Rees \Term{lemma}{引理}”（4.2.6）即得结论。
\end{proof}

\begin{SGAresult}{推论}{5.3.2}
\Term{delta-functor}{$\partial$ 函子}
\[
T=(T^i)_{i\in\mathbb N}:\operatorname{Hom}(\mathbb N^\circ,C)\longrightarrow\operatorname{Hom}(\mathbb N^\circ,D)
\]
诱导一个 $A$ \Term{linear}{线性}\Term{exact}{正合}\Term{delta-functor}{$\partial$ 函子}
\[
\mathscr E_C\longrightarrow\mathscr E_D,
\]
它与\Term{translation}{平移}交换；所以取商后得到一个新的 $A$ \Term{linear}{线性}\Term{exact}{正合}\Term{delta-functor}{$\partial$ 函子}，仍记作 $(T^i)_{i\in\mathbb N}$：
\[
\mathrm{AR}\text{-}J\text{-}\mathrm{adn}(C)\longrightarrow\mathrm{AR}\text{-}J\text{-}\mathrm{adn}(D).
\]
\end{SGAresult}


\SourcePage{255}
\SGAsection{附录}{Shih 定理}
本附录的目的是证明 5.3.1 的证明中所用\Term{form}{形式}的 Shih \Term{theorem}{定理}。由于这一\Term{form}{形式}比 EGA $0_{\mathrm{III}}$ 中的\Term{form}{形式}更一般，有必要给出一个由该处证明改写而成的直接证明。

\textbf{术语。} 设 $r_0$ 是一个 $\geq1$ 的整数。对每个整数 $r\geq r_0$，给定一个带有双\Term{degree}{次数}为 $(r,1-r)$ 的\Term{differential}{微分} $d_r$ 的\Term{bigraded-object}{双分次对象}
\[
E_r=(E_r^{pq})_{p\in\mathbb Z,\,q\in\mathbb Z},
\]
并对每个 $r$ 给定一个\Term{isomorphism}{同构}
\[
H(E_r)\xrightarrow{\sim}E_{r+1}.
\]
以下称这样的数据为一个\Term{spectral-filter}{谱滤}（filtre spectral）。\Term{spectral-filter}{谱滤}\Term{morphism}{态射}的概念按显然的方式定义。

\textbf{A.1. 回顾与 $\partial$ \Term{functor}{函子}相伴的\Term{spectral-sequence}{谱序列}。}
设 $C$ 是一个\Term{abelian-category}{阿贝尔范畴}，$X$ 是 $C$ 的一个\Term{object}{对象}，带有递减\Term{filtration}{滤过} $(F^pX)_{p\in\mathbb Z}$。对每个从 $C$ 到一个\Term{abelian-category}{阿贝尔范畴} $D$ 的\Term{exact}{正合}\Term{delta-functor}{$\partial$ 函子} $T$，所谓“\Term{exact-couple}{正合偶}”的方法（例如参见 Mac Lane，\textit{Homology}，第 336 页）给出一个\Term{spectral-sequence}{谱序列}
\[
E(X):E_1^{pq}=T^{p+q}(F^pX/F^{p+1}X)\Longrightarrow E^n=T^n(X),
\]
其中 $E^n$ 的\Term{filtration}{滤过}由下式定义：
\[
\begin{split}
F^p(E^n)&=\operatorname{Im}\bigl(T^n(F^pX)\longrightarrow T^n(X)\bigr)\\
&=\operatorname{Ker}\bigl(T^n(X)\longrightarrow T^n(X/F^pX)\bigr).
\end{split}
\]
此外，若 $X$ 的\Term{filtration}{滤过}是\Term{finite}{有限}的，则如此得到的\Term{spectral-sequence}{谱序列}是\Term{weak-convergence}{弱收敛}的（EGA $0_{\mathrm{III}}$ 11.1.3）。

\SourcePage{256}
\textbf{\Term{functorial}{函子性}。} 若 $f:X\to Y$ 是 $C$ 中的\Term{filtered-object}{滤过对象}的一个\Term{morphism}{态射}，则可\Term{functorial}{函子性}地构造一个\Term{spectral-sequence}{谱序列}\Term{morphism}{态射}
\[
E(f):E(X)\longrightarrow E(Y),
\]
它在 $E_1^{pq}(X)$（相应地，$E^n(X)$）上等于 $T^{p+q}(\operatorname{gr}^p f)$（相应地，$T^n(f)$）。
特别地，若 $u:X\to X$ 是 $X$ 的一个\Term{endomorphism}{自同态}，满足
\[
u(F^pX)\subset F^{p+h}(X),
\]
其中整数 $h$ 与 $p$ 无关，则 $E(u)$ 是与 $E(X)$ 相伴的\Term{spectral-filter}{谱滤}的一个双\Term{degree}{次数}为 $(h,-h)$ 的\Term{endomorphism}{自同态}，且与各\Term{differential}{微分}交换。

\textbf{A.2. \Term{inverse-system}{投射系}的情形。} 现在设 $X=(X_n)_{n\in\mathbb Z}$ 是一个\Term{abelian-category}{阿贝尔范畴} $C$ 的\Term{object}{对象}所构成的\Term{projective-system}{投射系}，满足
\[
X_i=0\qquad(i<0).
\]
若 $T$ 是从 $C$ 到一个\Term{abelian-category}{阿贝尔范畴} $D$ 的\Term{exact}{正合}\Term{delta-functor}{$\partial$ 函子}，则 $T$ 自然延拓为一个仍记作 $T$ 的\Term{exact}{正合}\Term{delta-functor}{$\partial$ 函子}
\[
\operatorname{Hom}(\mathbb Z^\circ,C)\longrightarrow\operatorname{Hom}(\mathbb Z^\circ,D).
\]
于是，赋予 $X$ 自然的\Term{projective-system}{投射系}\Term{filtration}{滤过}（EGA $0_{\mathrm{III}}$ 13.4.1），便得到一个\Term{spectral-sequence}{谱序列}
\[
E(X)=E(X_n)_{n\in\mathbb Z},
\]
其中 $E(X_n)$ 表示与 $T$ 以及带有诱导\Term{filtration}{滤过}的 $X_n$ 相伴的\Term{spectral-sequence}{谱序列}。

\textbf{$E(X)$ 的性质。}

(a) 对每个整数 $n$，当 $X$ 是\Term{strict-projective-system}{严格}的时，\Term{spectral-sequence}{谱序列} $E(X)$ 的\Term{abutment}{终项} $E^n(X)$ 上的\Term{filtration}{滤过}等于其自然的\Term{projective-system}{投射系}\Term{filtration}{滤过}。
事实上，以字母 $F$（同 A.1）表示前一种\Term{filtration}{滤过}，
\SourcePage{257}
以字母 $\Phi$ 表示后一种\Term{filtration}{滤过}，则对每个整数 $\alpha$ 都有
\[
\begin{split}
F^pE^n(X)_\alpha
&=\operatorname{Ker}\bigl(T^n(X_\alpha)\longrightarrow T^n(X_\alpha/F^pX_\alpha)\bigr)\\
&=\begin{cases}
0,&\alpha\leq p-1,\\
\operatorname{Ker}\bigl(T^n(X_\alpha)\longrightarrow T^n(X_{p-1})\bigr),&\alpha\geq p-1.
\end{cases}
\end{split}
\]
两种情形下，这正是 $\Phi^pE^n(X)$ 的表达式。
特别地，这表明当 $\beta\geq\alpha\geq p$ 时，\Term{transition-morphism}{转移态射}
\[
\begin{tikzcd}[column sep=3em]
E_\infty^{pq}(X_\beta)\arrow[r]\arrow[d,equal]&E_\infty^{pq}(X_\alpha)\arrow[d,equal]\\
\operatorname{gr}^p T^{p+q}(X_\beta)\arrow[r]&\operatorname{gr}^p T^{p+q}(X_\alpha)
\end{tikzcd}
\]
是一个\Term{monomorphism}{单态射}，这里使用了 EGA $0_{\mathrm{III}}$ 13.4.1.3。

(b) 固定 $\alpha\in\mathbb Z$，由 $i\geq\alpha+1$ 时 $F^i(X_\alpha)=0$，显然有（参看下图）
\[
\begin{aligned}
B_r^{pq}(X_\alpha)&=B_{r+1}^{pq}(X_\alpha)&&\text{当 }r>p,\\
Z_r^{pq}(X_\alpha)&=Z_{r+1}^{pq}(X_\alpha)&&\text{当 }r>\alpha-p.
\end{aligned}
\]
\[
\begin{tikzpicture}[scale=.8]
\draw (-.7,0)--(4.2,0);
\draw (0,-.35)--(0,3.4);
\draw (2.7,-.35)--(2.7,3.4);
\node[below left] at (0,0) {$0$};
\node[below right] at (2.7,0) {$\alpha$};
\fill (2.7,2.5) circle (1.2pt);
\node[above left] at (2.7,2.5) {$(p,q)$};
\draw[->] (2.1,2.95)--(3.8,1.95);
\draw[->] (2.1,2.95)--(2.1,1.95) node[midway,left] {$1-r$};
\draw[->] (2.1,1.95)--(3.8,1.95) node[midway,above] {$r$};
\end{tikzpicture}
\]

(c) 当 $X$ 满足\Term{ml-condition}{Mittag--Leffler 条件}时，各\Term{projective-system}{投射系} $\operatorname{gr}^p(X)$ 都是\Term{essentially-constant}{本质常值}的（EGA $0_{\mathrm{III}}$ 13.4.3）。由此立即得到，对于 $r<+\infty$，\Term{projective-system}{投射系} $E_r^{pq}(X)$ 都是\Term{essentially-constant}{本质常值}的。以
\[
\mathscr E_r^{pq}(X)
\]
表示它们的\Term{projective-limit}{投射极限}。将\Term{differential-operator}{微分算子}取\Term{limit}{极限}，便得到一个\Term{spectral-filter}{谱滤}
\[
\mathscr E_r^{pq}(X)\qquad(p,q\in\mathbb Z;\ r\geq1).
\]
与此\Term{spectral-filter}{谱滤}相伴的\Term{object}{对象}
\[
\mathfrak Z_r^{pq}(X),\qquad\mathfrak B_r^{pq}(X)\qquad(r<+\infty)
\]
\SourcePage{258}
恰是\Term{essentially-constant}{本质常值}\Term{projective-system}{投射系} $Z_r^{pq}(X)$ 与 $B_r^{pq}(X)$ 的\Term{limit}{极限}。
由 (b) 还得到
\[
B_\infty^{pq}(X)=B_{p+1}^{pq}(X),
\]
从而可以置
\[
\mathfrak B_\infty^{pq}(X)=\mathfrak B_{p+1}^{pq}(X).
\]
仍假定 $X$ 满足\Term{ml-condition}{$(\mathrm{ML})$}。由于对每个 $\alpha$，只要 $r>p$ 就有 $B_r^{pq}(X_\alpha)=B_{r+1}^{pq}(X_\alpha)$（参见 (b)），我们得到一列\Term{canonical}{典范}\Term{monomorphism}{单态射}：
\[
\cdots\subset\mathscr E_r^{pq}(X)\subset\cdots\subset\mathscr E_{p+2}^{pq}(X)\subset\mathscr E_{p+1}^{pq}(X).
\tag{S}
\]
以下假定 $X$ 是\Term{strict-projective-system}{严格}的。$Z_r^{pq}(X_n)$ 与 $B_r^{pq}(X_n)$ 的下列表达式
\[
\begin{split}
Z_r^{pq}(X_n)=\operatorname{Im}\bigl[&T^{p+q-1}(F^{p-r+1}X_n/F^pX_n)\\
&\xrightarrow{\partial}T^{p+q}(F^pX_n/F^{p+1}X_n)\bigr],
\end{split}
\]
\[
\begin{split}
B_r^{pq}(X_n)=\operatorname{Im}\bigl[&T^{p+q}(F^pX_n/F^{p+r}X_n)\\
&\longrightarrow T^{p+q}(F^pX_n/F^{p+1}X_n)\bigr]
\end{split}
\]
表明
\[
\left\{\begin{aligned}
Z_r^{pq}(X_m)&\xrightarrow{\sim}Z_r^{pq}(X_n)&&\text{当 }m\geq n\geq p,\\
B_r^{pq}(X_m)&\xrightarrow{\sim}B_r^{pq}(X_n)&&\text{当 }m\geq n\geq p+r-1
\end{aligned}\right.
\tag{$d_1$}
\]
（这些\Term{canonical-morphism}{典范态射}都是\Term{isomorphism}{同构}）。特别地，\Term{canonical-morphism}{典范态射}
\[
\mathscr E_r^{pq}(X)\longrightarrow E_r^{pq}(X_{p+r-1})
\]
是一个\Term{isomorphism}{同构}。
另一方面，由 (b) 得到
\SourcePage{259}
\[
\left\{\begin{aligned}
Z_\infty^{pq}(X_{p+r-1})&=Z_r^{pq}(X_{p+r-1})&&(r\geq1),\\
B_\infty^{pq}(X_{p+r-1})&=B_r^{pq}(X_{p+r-1})&&(r>p).
\end{aligned}\right.
\tag{$d_2$}
\]
这使我们能对每个 $r$ 定义一个\Term{canonical}{典范}\Term{epimorphism}{满态射}
\[
\theta_r^{pq}:\mathscr E_r^{pq}(X)\longrightarrow E_\infty^{pq}(X_{p+r-1}),
\]
当 $r>p$ 时，它是一个\Term{isomorphism}{同构}。此外，若整数 $a$ 满足
\[
\mathfrak B_r^{pq}(X)=\mathfrak B_a^{pq}(X)\qquad(r\geq a)
\]
（例如可取 $a=p+1$），从而有
\[
\cdots\lhook\joinrel\longrightarrow\mathscr E_r^{pq}(X)\lhook\joinrel\longrightarrow\cdots\lhook\joinrel\longrightarrow\mathscr E_{a+1}^{pq}(X)\longrightarrow\mathscr E_a^{pq}(X),
\tag{S bis}
\]
则对 $s\geq r\geq a$，下图交换：
\[
\begin{tikzcd}[column sep=3em]
\mathscr E_s^{pq}(X)\arrow[r,"\theta_s^{pq}"]\arrow[d,hook]&E_\infty^{pq}(X_{p+s-1})\arrow[d,"\mathrm{can}"]\\
\mathscr E_r^{pq}(X)\arrow[r,"\theta_r^{pq}"]&E_\infty^{pq}(X_{p+r-1}).
\end{tikzcd}
\tag{$d_3$}
\]

\textbf{A.3. Shih \Term{theorem}{定理}的一个变式。} 假定给定一个环 $S$，带有由\Term{ideal}{理想} $F^i(S)$（$i\geq0$）组成的\Term{filtration}{滤过}，并且 $S$ 在 $X$ 上的\Term{action}{作用}与其\Term{projective-system}{投射系}\Term{filtration}{滤过}相容。
由 A.1 的一般论述，对每个整数 $i$，$\operatorname{gr}^i(S)$ 都通过双\Term{degree}{次数}为 $(i,-i)$ 的\Term{morphism}{态射}\Term{action}{作用}于\Term{bigraded-object}{双分次对象} $E_r(X)$。特别地，对每个整数 $n$，
\[
E_r^{(n)}(X)=(E_r^{pq}(X))_{p+q=n}
\]
\Term{canonical}{典范}地具有\Term{graded-module}{分次 $(\operatorname{gr}^{\bullet}(S),D)$ 模}的结构。同样，取\Term{limit}{极限}可得：
\SourcePage{260}
\[
E_r^{(n)}(X)=(E_r^{pq}(X))_{p+q=n}
\]
\Term{canonical}{典范}地具有\Term{graded-module}{分次 $(\operatorname{gr}^{\bullet}(S),D)$ 模}的结构。同样，取\Term{limit}{极限}可得
\[
\mathscr E_r^{(n)}(X)=(\mathscr E_r^{pq}(X))_{p+q=n}.
\]
环 $S$ 在 $X$ 上的\Term{action}{作用}与其\Term{projective-system}{投射系}\Term{filtration}{滤过}相容，因而也对每个整数 $n$ \Term{action}{作用}于 $T^n(X)$。所以\Term{strict-graded}{严格分次对象} $\operatorname{grs}(X)$ \Term{canonical}{典范}地具有\Term{graded-module}{分次 $(\operatorname{gr}^{\bullet}(S),C)$ 模}的结构；同样，当 $T^n(X)$ 满足\Term{ml-condition}{$(\mathrm{ML})$} 时，\Term{strict-graded}{严格分次对象} $\operatorname{grs}^{\bullet}T^n(X)$ \Term{canonical}{典范}地具有\Term{graded-module}{分次 $(\operatorname{gr}^{\bullet}(S),D)$ 模}的结构。

\begin{SGAresult}{\Term{theorem}{定理}}{A.3}
假定对某个整数 $n$，\Term{graded-module}{分次 $(\operatorname{gr}^{\bullet}(S),D)$ 模}
\[
\mathscr E_1^{(n)}(X)=T^n(\operatorname{grs}^{\bullet}(X)),\qquad
\mathscr E_1^{(n-1)}(X)=T^{n-1}(\operatorname{grs}^{\bullet}(X))
\]
都是\Term{noetherian}{诺特}的。则：

(i) $T^n(X)$ 满足\Term{mlar-condition}{$(\mathrm{MLAR})$}；

(ii) $\operatorname{grs}^{\bullet}T^n(X)$ 是一个\Term{noetherian}{诺特}的 $(\operatorname{gr}^{\bullet}(S),D)$ 模。
\end{SGAresult}
\begin{proof}
(i) 对每对整数 $(m,r)$，
\[
\mathfrak B_r^{(m)}(X)=\bigl(\mathfrak B_r^{p,m-p}(X)\bigr)_{p\in\mathbb N}
\]
都是 $\mathscr E_1^{(m)}(X)$ 的一个\Term{graded-module}{分次 $(\operatorname{gr}^{\bullet}(S),D)$ 子模}。特别地，对每个整数 $m$，各 $\mathfrak B_r^{(m)}(X)$ 构成 $\mathscr E_1^{(m)}(X)$ 的\Term{graded-module}{分次 $(\operatorname{gr}^{\bullet}(S),D)$ 子模}的一个以 $r$ 为指标的递增列。利用 $\mathscr E_1^{(n)}(X)$ 与 $\mathscr E_1^{(n-1)}(X)$ 都是\Term{noetherian}{诺特}的\Term{graded-module}{分次 $(\operatorname{gr}^{\bullet}(S),D)$ 模}这一事实，可知两列
\SourcePage{261}
\[
\bigl[\mathfrak B_r^{(n)}(X)\bigr]_{r\in\mathbb N},\qquad
\bigl[\mathfrak B_r^{(n-1)}(X)\bigr]_{r\in\mathbb N}
\]
都最终\Term{invariant}{不变}。

因此存在整数 $a\geq1$，使得对于每个满足 $p+q=n$ 或 $n-1$ 的 $(p,q)$，都有
\[
\mathfrak B_r^{pq}(X)=\mathfrak B_a^{pq}(X)\qquad(r\geq a).
\tag{A.3.1}
\]
于是，当 $r\geq a+1$ 时，以 $\mathscr E_r^{pq}(X)$（$p+q=n$）为源或为靶的\Term{arrow}{箭头} $d_r$ 均为零。换言之，由 $p+q=n$ 时的 A.3.1 得到的\Term{monomorphism}{单态射}
\[
\mathscr E_s^{pq}(X)\lhook\joinrel\longrightarrow\mathscr E_r^{pq}(X)\qquad(s\geq r\geq a+1)
\]
实际上都是\Term{isomorphism}{同构}。对于 $s\geq r\geq a+1$，$p+q=n$ 时图 $(d_3)$ 的\Term{commutativity}{交换性}，即
\[
\begin{tikzcd}[column sep=3em]
\mathscr E_s^{pq}(X)\arrow[r,"\theta_s^{pq}",two heads]\arrow[d,"\sim"']&E_\infty^{pq}(X_{p+s-1})\arrow[d]\\
\mathscr E_r^{pq}(X)\arrow[r,"\theta_r^{pq}",two heads]&E_\infty^{pq}(X_{p+r-1}),
\end{tikzcd}
\]
表明\Term{transition-morphism}{转移态射}
\[
E_\infty^{pq}(X_{p+s-1})\longrightarrow E_\infty^{pq}(X_{p+r-1})
\]
是一个\Term{epimorphism}{满态射}；而它总是一个\Term{monomorphism}{单态射}（参见 (a)），因而是一个\Term{isomorphism}{同构}。换言之，\Term{projective-system}{投射系}
\[
r\longmapsto\left[\bigoplus_p\operatorname{gr}^pT^n(X_{p+r})\right]_{r\geq1}
\]
从 $r=a$ 起为常值，故由 4.1.3，$T^n(X)$ 满足\Term{mlar-condition}{$(\mathrm{MLAR})$}。

(ii) 以 $\mathscr E_\infty^{(n)}(X)$、$\mathfrak Z_\infty^{(n)}(X)$、$\mathfrak B_\infty^{(n)}(X)$ 分别表示所考虑的各个以 $r$ 为指标的\Term{projective-system}{投射系}的“\Term{limit}{极限}”，则有
\[
\mathscr E_\infty^{(n)}(X)=\mathfrak Z_\infty^{(n)}(X)/\mathfrak B_\infty^{(n)}(X).
\]
\SourcePage{262}
这是一个\Term{noetherian}{诺特}的 $(\operatorname{gr}^{\bullet}(S),D)$ 模，因为 $\mathfrak Z_\infty^{(n)}(X)$ 包含于 $\mathscr E_1^{(n)}(X)$。另一方面，由 A.2 (a) 得到
\[
\mathscr E_\infty^{(n)}(X)=\operatorname{grs}^{\bullet}T^n(X).
\]
\end{proof}